\section{Associative Memory in Biological and Artificial Systems}

Associative memory refers to the ability of a system to store and retrieve information based on relationships between patterns. In artificial intelligence (AI), associative memory models are designed to recover stored information from partial or incomplete inputs by learning associations between representations. Early approaches, such as Hopfield Networks, modelled memories as stable states within interconnected networks, where the structure of connections influenced memory storage and retrieval \cite{hopfield1982}. Modern Hopfield Networks (MHNs) build upon this foundation by introducing high-dimensional representations and attention-based retrieval mechanisms, allowing associative memory models to achieve greater scalability and improved retrieval capabilities \cite{ramsauer2021hopfield}.

The concept of associative memory is also observed in biological neural systems, where it supports processes such as learning, recognition, and recall. In the human brain, memories are represented through distributed patterns of neural activity across interconnected networks rather than being stored in individual neurons \cite{kandel2013principles}. The connections between neurons are formed through synapses, whose strengths can change through synaptic plasticity, allowing the brain to adapt based on experience. One proposed mechanism for this adaptation is Hebbian learning, where repeated co-activation between neurons strengthens their connections and contributes to the formation of associations between patterns \cite{hebb1949organization}. 